LANCE est un agent fondé sur un modèle de langage pour conduire et évaluer des
audits de réseaux IoT/OT. La mission principale du stage a consisté à concevoir
une architecture HMoE agentique~: un modèle de base quantifié est partagé entre
des adaptateurs QLoRA spécialisés pour l'organisation, la reconnaissance,
l'analyse de vulnérabilités et l'exploitation. Une chaîne de préparation
transforme les traces en exemples validés, tandis qu'un routeur déterministe
sélectionne l'expert associé à chaque phase derrière une API commune. Le
pipeline encadre les transitions, les outils, les sorties et les preuves. En
parallèle, le harness expérimental déploie des scénarios reproductibles et mesure
le comportement observé pendant les campagnes. L'évaluation associe des
métriques déterministes à un jugement sémantique encadré, sans confier le calcul
final au modèle. Des mécanismes de maîtrise du contexte complètent
l'architecture. Les premiers résultats
mettent en évidence un écart entre constats produits et chemins d'attaque
entièrement prouvés. Ils doivent encore être confirmés par des campagnes répétées
comparant les experts spécialisés à une référence généraliste dans des
conditions identiques.

\vspace{1em}
\noindent\textbf{Mots-clés~:} agent LLM, cybersécurité IoT/OT, HMoE.
